\newpage\section{对称性}
一个\textbf{对称变换}改变的是观察方式，而非物理系统本身，因而不改变任何可能的实验结果。
从\textbf{被动观点}看，物理系统并未发生任何变化，改变的只是描述它的参考框架（或观测者）。设观测者 $O$ 用射线 $\mathscr{R},\,\mathscr{R}_1,\,\mathscr{R}_2,\dots$ 描述同一系统的各个态，而等价观测者 $O'$（例如相对于 $O$ 做了旋转或平移）用射线 $\mathscr{R}',\,\mathscr{R}'_1,\,\mathscr{R}'_2,\dots$ 来描述\textbf{同一个}系统的对应态。所谓"对称"，就是要求两位观测者计算的跃迁几率完全一致：
\begin{equation}
    P(\mathscr{R} \to \mathscr{R}_n) = P(\mathscr{R}' \to \mathscr{R}'_n)
\end{equation}

换言之：
\begin{itemize}
    \item \textbf{不变的}：物理系统本身、实验结果（几率）。
    \item \textbf{改变的}：观测者（参考框架），以及随之而来的态的数学描述（射线 $\mathscr{R}\mapsto\mathscr{R}'$）。
\end{itemize}
这正是\textbf{被动观点}的核心：对称变换并非"动了系统"，而是"换了观测者"，物理规律（几率）在这一更换下保持不变。
\par\vspace{2ex}
    \hrule
\vspace{2ex}
\footnote{Wigner 的证明省略了一些步骤。一个更加完整的证明参考S.Weinberg书}
    Wigner在20世纪30年代早期所证明的一个重要定理告诉我们,对于任意这样的射线变换 $\mathscr{R} \to \mathscr{R}'$，我们可以定义 Hilbert 空间上的一个算符 $U$，使得如果 $|\psi\rangle$ 在射线 $\mathscr{R}$ 中，那么 $U|\psi\rangle$ 在射线 $\mathscr{R}'$ 中，则 $U$ 要么是么正且线性的：
\begin{align}
    \langle U\varphi| U\psi \rangle &= \langle \varphi| \psi \rangle \;, \\
    U(\alpha|\varphi\rangle + \beta|\psi\rangle) &= \alpha U|\varphi\rangle + \beta U|\psi\rangle \;.
\end{align}
要么是反么正且反线性的：
\begin{align}
    \langle U\varphi| U\psi \rangle &= \langle \varphi| \psi \rangle^* \;, \label{eq:anti_unitary_prod}\\
    U(\alpha|\varphi\rangle + \beta|\psi\rangle) &= \alpha^* U|\varphi\rangle + \beta^* U|\psi\rangle \;.
\end{align}
这里么正或反么正的条件均采取以下形式\footnote{反线性算符 $A$ 的共轭定义为
\begin{equation}
    \langle \varphi| \mathscr{O}^\dagger | \psi \rangle = \langle A\mathscr{O}\varphi| \psi \rangle^* = \langle \psi | \mathscr{O}\varphi\rangle \;.
\end{equation}}
\begin{equation}
    U^\dagger = U^{-1} \;.
\end{equation}
 平庸的对称变换 $\mathscr{R} \to \mathscr{R}$，是单位算符 $U=\mathbf{1}$ 显然，单位算符 $U=\mathbf{1}$ 是线性且么正的。由于\textbf{连续对称性}可以通过调节参数连续地变回单位算符，因此它们必须由线性么正算符表示。而反线性反么正算符仅用于涉及时间反演的非连续变换。
\vspace{2ex}\par
 特别的，一个无限小的接近于平庸的对称变换可以表示为一个线性么正算符：$$U = \mathbf{1} + i\epsilon t$$其中$\epsilon$ 是一个实的无限小参量。由于 $U$ 是么正且线性的，$t$ 必须厄米且线性；所以它可能是可观测量，确实，大多数（也许是全部）物理的可观测量从对称变换中诞生。
\subsection{对称群与投影表示}\label{sec:projective}

对称变换构成一个群：两次连续变换 $T_1, T_2$ 等价于另一次变换 $T_3 = T_2 T_1$；存在逆变换 $T^{-1}$ 和单位变换 $\mathbf{1}$。数学表述是：
\begin{tcolorbox}[colback=white!5,colframe=black!45]
对称变换的集合 $\mathcal{G}$ 在\textbf{群乘法}（连续变换）下构成一个\textbf{群}。
\textbf{乘法定义}：若 $T_1$ 将射线 $\mathscr{R}$ 变换为 $\mathscr{R}'$，$T_2$ 将 $\mathscr{R}'$ 变换为 $\mathscr{R}''$，则乘积 $T_2 T_1$ 定义为将 $\mathscr{R}$ 直接变换为 $\mathscr{R}''$ 的单一变换。
该集合满足以下群公理：
\begin{itemize}
    \item \textbf{封闭性}：任意两个变换的乘积仍属于该集合，即 $\forall T_1, T_2 \in \mathcal{G} \implies T_2 T_1 \in \mathcal{G}$；
    \item \textbf{结合律 }：变换的顺序满足 $(T_3 T_2) T_1 = T_3 (T_2 T_1)$；
    \item \textbf{单位元 }：存在恒等变换 $\mathbf{1} \in \mathcal{G}$，使得 $\forall T \in \mathcal{G}, T \mathbf{1} = \mathbf{1} T = T$；
    \item \textbf{逆元 }：对任意变换 $T$，存在逆变换 $T^{-1}$ 使得 $T T^{-1} = T^{-1} T = \mathbf{1}$。
\end{itemize}
\end{tcolorbox}
由于量子力学中的物理态由希尔伯特空间中的\textbf{射线 (Ray)} 而非矢量唯一确定（即 $|\psi\rangle$ 与 $e^{i\alpha}|\psi\rangle$ 描述同一物理态），对称性变换算符 $U(T)$ 必须满足 Wigner 定理，即 $U(T)$ 是幺正或反幺正的。然而，这一性质使得群的乘法规则在希尔伯特空间中表现得更为复杂。当两个对称操作连续作用时，其结果只需在“射线意义上”与群乘法一致。这意味着算符的乘积关系可以相差一个相位因子：
\begin{equation}
    U(T_2)U(T_1)|\psi_n\rangle = e^{i\varphi_n(T_2, T_1)} U(T_2 T_1)|\psi_n\rangle \;.
\end{equation}
乍看之下，相位 $\varphi_n$ 似乎可以依赖于具体的态 $|\psi_n\rangle$。但我们有如下重要结论：
\begin{tcolorbox}[colback=white!5,colframe=black!45]
        如果 $U(T)$ 是线性（或反线性）算符，则上述相位因子 $\varphi(T_2, T_1)$ 必须是\textbf{全局的}，即与具体状态 $|\psi\rangle$ 无关。
\end{tcolorbox}
\noindent
\\
\textbf{证明：}
利用量子力学的态叠加原理。考虑两个线性无关的态 $|\psi_A\rangle$ 和 $|\psi_B\rangle$，以及它们的叠加态 $|\psi_{AB}\rangle = |\psi_A\rangle + |\psi_B\rangle$。

一方面，直接作用于叠加态：
\begin{equation}
    U(T_2)U(T_1) |\psi_{AB}\rangle = e^{i\varphi_{AB}} U(T_2 T_1) (|\psi_A\rangle + |\psi_B\rangle) \;. \label{eq:phase_proof_1}
\end{equation}
另一方面，利用算符 $U$ 的线性（分配律），先分别作用再叠加：
\begin{equation}
    \begin{aligned}
        U(T_2)U(T_1) (|\psi_A\rangle + |\psi_B\rangle) &= U(T_2)U(T_1)|\psi_A\rangle + U(T_2)U(T_1)|\psi_B\rangle \\
        &= e^{i\varphi_A} U(T_2 T_1)|\psi_A\rangle + e^{i\varphi_B} U(T_2 T_1)|\psi_B\rangle \;.
    \end{aligned} \label{eq:phase_proof_2}
\end{equation}
比较式 \eqref{eq:phase_proof_1} 和 \eqref{eq:phase_proof_2} 的右边。由于 $|\psi_A\rangle$ 和 $|\psi_B\rangle$ 是线性无关的，变换后的态 $U(T_2 T_1)|\psi_A\rangle$ 和 $U(T_2 T_1)|\psi_B\rangle$ 依然保持线性无关（因为 $U$ 是可逆的幺正/反幺正算符）。

因此，要使等式成立，对应基底前的系数必须严格相等：
\begin{equation}
    e^{i\varphi_{AB}} = e^{i\varphi_A} \quad \text{且} \quad e^{i\varphi_{AB}} = e^{i\varphi_B} \;.
\end{equation}
这直接导致 $\varphi_A = \varphi_B = \varphi_{AB} \equiv \varphi(T_2, T_1)$。所以，对称群在希尔伯特空间上的实现遵循如下乘积规则：
\begin{equation}
    U(T_2)U(T_1) = e^{i\varphi(T_2, T_1)} U(T_2 T_1) \;. \label{eq:projective_rep}
\end{equation}
这种包含相位因子的表示被称为群的\textbf{投影表示 }。
\begin{itemize}
    \item 若对于所有 $T_1, T_2$ 均可选择相位使得 $\varphi=0$，则退化为\textbf{普通表示}。
    \item 在特定拓扑条件下（如庞加莱群），可以通过重新定义算符相位或考虑群的万有覆盖群，将投影表示提升为普通表示。
\end{itemize}
\topictitle{***}
在上面的讨论中，我们已经证明了对称群在Hilbert空间上的实现必然是\textbf{投影表示}：
\begin{equation}
    U(T_2)\,U(T_1) = e^{i\varphi(T_2,T_1)}\,U(T_2 T_1)\,.
\end{equation}
一个自然的问题是：这个相位因子$\varphi(T_2,T_1)$能否通过重新定义算符$U(T)\to e^{i\alpha(T)}U(T)$来消除？
如果能全部消除，投影表示就退化为普通表示，半整数自旋将不被允许。


答案取决于对称群的\textbf{拓扑结构}。
可以证明\footnote{%
    严格的数学表述涉及群上同调$H^2(G,\,U(1))$。
    对于单连通的李群，所有投影表示都可以提升为普通表示；
    对于非单连通的李群，存在不可消除的相位因子，
    其完整分类由万有覆盖群的表示给出。
    参见Weinberg, \textit{QFT} Vol.\,I, \S2.7。
}：
\begin{itemize}
    \item 若$G$是\textbf{单连通}的（$\pi_1(G)=0$），
    则所有投影表示的相位因子都可以消除——只需考虑普通表示。
    \item 若$G$\textbf{不是}单连通的（$\pi_1(G)\neq 0$），
    则存在\emph{不可消除}的相位因子。
    此时必须将$G$替换为其\textbf{万有覆盖群}$\widetilde{G}$，
    用$\widetilde{G}$的普通表示来完整分类$G$的全部投影表示。
\end{itemize}

因此，要判断旋转群$SO(3)$是否允许半整数自旋（即费米子是否在数学上合法），
\textbf{关键问题}归结为：

\begin{center}
\fbox{\parbox{0.75\textwidth}{\centering
$SO(3)$是单连通的吗？\\[4pt]
如果不是，它的万有覆盖群$\widetilde{SO(3)}$是什么？\\
$\widetilde{SO(3)}$是否拥有$SO(3)$本身不具有的额外表示？
}}
\end{center}

\noindent 下面我们通过具体的几何构造来回答这些问题。

\subsubsection*{$SO(3)$的参数空间}

三维空间中的任意旋转可以用一个旋转轴$\hat{\boldsymbol{n}}$和一个旋转角$\theta\in[0,\pi]$来指定。
我们可以将这对数据编码为一个三维矢量：
\begin{equation}
    \boldsymbol{v} = \theta\,\hat{\boldsymbol{n}}\,,\qquad |\boldsymbol{v}|\leq\pi\,.
\end{equation}
这样，$SO(3)$的每个群元素对应半径为$\pi$的实心球$B^3$中的一个点，
球心$\boldsymbol{v}=\mathbf{0}$对应恒等变换$\mathbf{1}$。

但是存在一个关键的\textbf{等价关系}：
绕$\hat{\boldsymbol{n}}$旋转$\pi$和绕$-\hat{\boldsymbol{n}}$旋转$\pi$是\emph{同一个}旋转%
\footnote{%
绕$\hat{\boldsymbol{n}}$旋转$\pi$的效果是将垂直于$\hat{\boldsymbol{n}}$的分量反号，
这与绕$-\hat{\boldsymbol{n}}$旋转$\pi$完全相同。}。
因此球面$|\boldsymbol{v}|=\pi$上的\textbf{对径点}必须被认同：
\begin{equation}
    \pi\,\hat{\boldsymbol{n}} \;\sim\; -\pi\,\hat{\boldsymbol{n}}\,.
\end{equation}
这意味着$SO(3)$的参数空间不是普通的实心球，而是对径点被粘合后的空间——即实射影空间：
\begin{equation}
    SO(3) \;\cong\; B^3\!/\!\sim\;\; = \;\mathbb{R}\mathrm{P}^3\,.
\end{equation}

\subsubsection*{可缩路径与不可缩路径}

现在考虑$SO(3)$中从恒等元$\mathbf{1}$出发并回到$\mathbf{1}$的闭合路径，
判断它们能否被连续收缩为一个点。

\begin{enumerate}
    \item \textbf{可缩路径}（图~\ref{fig:SO3-paths}左）。
    任何完全位于球\emph{内部}（不触及边界）的闭合环路，
    都可以连续地缩小为球心处的一个点。
    这类路径对应$SO(3)$中拓扑平凡的闭合回路。

    \item \textbf{不可缩路径}（图~\ref{fig:SO3-paths}右）。
    考虑如下路径：从球心$\mathbf{1}$出发，沿某个方向走到球面上的点$P$（旋转角$\theta=\pi$）；
    由于对径认同，$P$和对径点$\bar{P}$是同一个群元素，
    所以路径从$\bar{P}$``重新出现''，再走回球心$\mathbf{1}$。
    
    这是一条\textbf{闭合}路径（起点和终点都是$\mathbf{1}$），
    但它\emph{穿越}了被认同的边界一次。
    可以证明%
    \footnote{严格证明可参考Nakahara, \textit{Geometry, Topology and Physics}, \S4.7。}%
    这条路径\textbf{不能}被连续收缩为一个点——
    它``钩住''了$\mathbb{R}\mathrm{P}^3$中不可收缩的拓扑结构。
\end{enumerate}

然而，如果我们将不可缩路径走\emph{两遍}%
（即对应旋转角从$0$到$2\pi$再到$4\pi$），
闭合路径穿越边界\emph{两次}——这条路径\emph{可以}被连续收缩为一个点%
\footnote{%
直观的理解：两次穿越可以通过连续形变``从两侧同时收回''而消除，
类似于将绕柱子缠绕两圈的绳子解开。}。
因此$SO(3)$的基本群恰好是：
\begin{equation}\label{eq:pi1-SO3}
    \boxed{\pi_1\bigl(SO(3)\bigr) = \mathbb{Z}_2}
\end{equation}
$SO(3)$是\textbf{双连通}的，不是单连通的。

% ====== 图 ======
\begin{figure}[htbp]
\centering
\begin{tikzpicture}[>=stealth, scale=1.0]

% ====== 左图：可缩路径 ======
\begin{scope}[shift={(-4.2,0)}]
  % 实心球截面
  \shade[ball color=gray!8, opacity=0.4] (0,0) circle (2.5);
  \draw[gray!60, semithick] (0,0) circle (2.5);
  
  % 对径认同参考线
  \draw[gray!40, thin, densely dotted] (2.5,0) -- (-2.5,0);
  \draw[gray!40, thin, densely dotted] (0,2.5) -- (0,-2.5);
  
  % 可缩的闭合路径
  \draw[blue!70!black, line width=2pt,
    decoration={markings,
      mark=at position 0.15 with {\arrow[scale=1.4]{stealth}},
      mark=at position 0.55 with {\arrow[scale=1.4]{stealth}},
      mark=at position 0.85 with {\arrow[scale=1.4]{stealth}}
    }, postaction={decorate}]
    (0,0) 
    .. controls (0.3,0.8) and (1.0,1.2) .. (1.2,0.6)
    .. controls (1.4,-0.2) and (0.8,-0.8) .. (0.2,-0.5)
    .. controls (-0.4,-0.2) and (-0.2,0.3) .. (0,0);
  
  % 收缩暗示
  \draw[red!60, ->, thick] (0.9,0.3) -- (0.4,0.1);
  \node[red!60, font=\small] at (1.5,0.3) {可缩};
  
  % 球心
  \fill[black] (0,0) circle (2.5pt);
  \node[below left, font=\bfseries] at (0,0) {$\mathbf{1}$};
  
  \node[below, font=\small\bfseries] at (0,-3.0) {可缩路径};
\end{scope}

% ====== 右图：不可缩路径 ======
\begin{scope}[shift={(4.2,0)}]
  % 实心球截面
  \shade[ball color=gray!8, opacity=0.4] (0,0) circle (2.5);
  \draw[gray!60, semithick] (0,0) circle (2.5);
  
  % 对径认同参考线
  \draw[gray!40, thin, densely dotted] (2.5,0) -- (-2.5,0);
  \draw[gray!40, thin, densely dotted] (0,2.5) -- (0,-2.5);
  
  % 对径点坐标
  \pgfmathsetmacro{\Px}{2.5*cos(50)}
  \pgfmathsetmacro{\Py}{2.5*sin(50)}
  
  % 段1：球心 → 边界点 P
  \draw[blue!70!black, line width=2pt,
    decoration={markings,
      mark=at position 0.35 with {\arrow[scale=1.4]{stealth}},
      mark=at position 0.78 with {\arrow[scale=1.4]{stealth}}
    }, postaction={decorate}]
    (0,0) .. controls (0.4,0.6) and (0.8,1.3) .. ({\Px},{\Py});
  
  % 段2：对径点 P̄ → 球心
  \draw[orange!80!black, line width=2pt,
    decoration={markings,
      mark=at position 0.3 with {\arrow[scale=1.4]{stealth}},
      mark=at position 0.75 with {\arrow[scale=1.4]{stealth}}
    }, postaction={decorate}]
    ({-\Px},{-\Py}) .. controls (-0.8,-1.3) and (-0.4,-0.6) .. (0,0);
  
  % 对径认同标记
  \draw[red!60!black, thick, densely dashed, <->] 
    ({\Px+0.25},{\Py+0.2}) 
    .. controls (3.2,1.0) and (3.2,-1.0) ..
    ({-\Px-0.25},{-\Py-0.2});
  \node[red!60!black, font=\small, right] at (3.3,0) {$P\sim\bar{P}$};
  
  % 边界点标记
  \fill[blue!70!black] ({\Px},{\Py}) circle (3pt);
  \node[above right, font=\small\bfseries, blue!70!black] at ({\Px},{\Py}) {$P$};
  
  \fill[orange!80!black] ({-\Px},{-\Py}) circle (3pt);
  \node[below left, font=\small\bfseries, orange!80!black] at ({-\Px},{-\Py}) {$\bar{P}$};
  
  % 球心
  \fill[black] (0,0) circle (2.5pt);
  \node[below right, font=\bfseries] at (0.05,-0.05) {$\mathbf{1}$};
  
  \node[below, font=\small\bfseries] at (0,-3.0) {不可缩路径};
\end{scope}

\end{tikzpicture}
\caption{$SO(3)$的参数空间$B^3/\mathbb{Z}_2 = \mathbb{R}\mathrm{P}^3$的截面示意图。
圆盘表示半径为$\pi$的实心球的一个截面，球心为恒等元$\mathbf{1}$，
边界上的对径点被认同（$P\sim\bar{P}$）。
\textbf{左}：完全位于球内部的闭合路径，可以连续收缩为球心处的一个点。
\textbf{右}：从$\mathbf{1}$出发到达边界点$P$（\textcolor{blue!70!black}{蓝色}段），
由于$P$和$\bar{P}$被认同，路径从$\bar{P}$``重新出现''并返回$\mathbf{1}$（\textcolor{orange!80!black}{橙色}段）。
这条闭合路径穿越了被认同的边界一次，拓扑上不可收缩。}
\label{fig:SO3-paths}
\end{figure}

\subsubsection*{物理后果：费米子的合法性}

现在回到投影表示的问题。$SO(3)$不是单连通的，
因此存在不可消除的相位因子$\varphi(T_2,T_1)$。
其万有覆盖群是$SU(2)$（群流形为单连通的三维超球面$S^3$），
覆盖映射$\rho: SU(2)\to SO(3)$是$2:1$的：
$SU(2)$中的$+\mathbf{I}$和$-\mathbf{I}$对应$SO(3)$中同一个旋转。

上述不可缩路径在$SU(2)$中被``展开''为从$+\mathbf{I}$到$-\mathbf{I}$的\emph{开}路径——
这正是旋量在$2\pi$旋转下获得负号的几何起源。
$SU(2)$除了拥有$SO(3)$的全部整数自旋表示（$j=0,1,2,\ldots$）之外，
还额外拥有\textbf{半整数自旋表示}（$j=1/2,3/2,\ldots$）。
在半整数自旋表示中：
\begin{equation}
    U(2\pi) = -\mathbf{I}\,.
\end{equation}
由于量子力学中$|\psi\rangle$和$-|\psi\rangle$描述同一个物理态（射线等价性），
这个负号在物理上是不可观测的——半整数自旋表示在物理上是\emph{完全合法}的。

\begin{tcolorbox}[colback=white!5, colframe=black!45,]
\textbf{费米子的存在}不是一个需要额外假设的经验事实，
而是以下两个基本原理的\textbf{不可避免的数学推论}：
\begin{enumerate}
    \item 量子力学的射线等价性（态由射线而非矢量确定）；
    \item 空间旋转群$SO(3)$的拓扑结构（$\pi_1 = \mathbb{Z}_2$，双连通）。
\end{enumerate}
同样的逻辑在相对论中再次上演：
固有正时Lorentz群$SO^+(1,3)$也是双连通的（$\pi_1=\mathbb{Z}_2$），
其万有覆盖群是$SL(2,\mathbb{C})$。
为了描述自然界中所有可能的基本粒子，
我们必须使用覆盖群$SL(2,\mathbb{C})$的表示，而非$SO^+(1,3)$自身的表示。
\end{tcolorbox}


\subsection{李群与李代数}
物理中最重要的对称群是\textbf{连通李群} (Connected Lie Group)，即可以通过连续参数 $\theta^a$ 描述、且能连续变回单位元的群。这个群的乘法规则是：
\begin{equation}
    T(\xi)T(\theta)=T\left\{g(\xi,\theta)\right\}\footnote{这里，$g^a(\xi,\theta)$是$\xi$和$\theta$的函数,由于李群在有限领域内都可以有限变换到单位元令$\xi^a=0$或$\theta^a=0$是单位元的坐标,必须有
    $$g^a(0,\theta)=g^a(\xi,0)=\theta^a\text{或}\xi^a$$且展开到二阶为：
    $$g^a(\xi,\theta)=\xi^a+\theta^a+f^a_{bc}\xi^b\theta^c+\ldots$$}
\end{equation}
对于李群，我们关注其在单位元附近的性质。将幺正算符 $U(T(\theta))$ 在 $\theta=0$ 处展开：
\begin{equation}
    U\left\{T(\theta)\right\} = \mathbf{1} + i\theta^a t_a + \frac{1}{2}\theta^b \theta^c t_{bc} + \dots \;
\end{equation}
对于两次变换,群表示变为：
\begin{equation}
    U\left\{ T(\xi)\right\}U\left\{ T(\theta)\right\}=U\left\{ T(g^a(0,\theta))\right\}
\end{equation}
上式：\[\begin{aligned}L.H.S=
    U\left\{ T(\theta)\right\}U\left\{ T(\xi)\right\}
    &=
   (\mathbf{1} + i\theta^a t_a + \frac{1}{2}\theta^b \theta^c t_{bc} + \dots \;)(\mathbf{1} + i\xi^a t_a + \frac{1}{2}\xi^b \xi^c t_{bc} + \dots \;)\\
    &=
    \mathbf{1}+i(\xi^a+\theta^a)t_a
    +\frac{1}{2}t_{bc}(\theta^b \theta^c+\xi^b \xi^c)-\xi^b\theta^c t_bt_c+\dots \;
\end{aligned}\]
\[\begin{aligned}R.H.S=
    U\left\{ T(g^a(0,\theta))\right\}
    &=U\left\{T(\xi^a+\theta^a+f^a_{bc}\xi^b\theta^c+\ldots)\right\}\\
    &=
    \mathbf{1}+
    i(\xi^a+\theta^a+f^a_{bc}\xi^b\theta^c+\ldots)t_a+
    \frac{1}{2}(\xi^b+\theta^b+\ldots)(\xi^c+\theta^c+\ldots)t_{bc}
    \\&=
    \mathbf{1}+i(\xi^a+\theta^a)t_a+if^a_{bc}\xi^b\theta^c t_a
    +\frac{1}{2}(\xi^b\xi^c+\xi^b\theta^c+\theta^c\xi^b+\theta^b\theta^c)t_{bc}
\end{aligned}\]
比较展开式中 $\bar{\theta}\theta$ 的交叉项系数，我们得到二阶项与一阶项的约束：
\begin{equation}
    \begin{aligned}
    &\xi^b\theta^c(t_bt_c+if^a_{bc} t_a+t_{bc})=0\\
    &\xi^b\theta^c(t_ct_b+if^a_{cb} t_a+t_{cb})=0\qquad\text{交换b和c的指标,且$t_{cb}$是对称的}
\end{aligned}
\end{equation}两式相减，得到导出了生成元必须满足的\textbf{对易关系}——即\textbf{李代数}：
\begin{equation}
    [t_b,t_c]=i(f^a_{cb}-f^a_{bc})t_a=iC^a_{bc}t_a
\end{equation}
\topictitle{阿贝尔群}
若群的变换可以任意交换顺序（如平移群），即为\textbf{阿贝尔群}。此时所有生成元对易，$C^a_{\; bc} = 0$，算符乘积的相位因子通常可以取为 1。
在此情况下，生成元相互独立，算符直接写为指数形式：
\begin{equation}
    U(T(\theta)) = \exp(i\theta^a t_a) \;.
\end{equation}


